% ============================================================
% SECTION 4 OPENING + 4.1 SETUP. Written 2026-07-26 (late).
%
% Part of the restructure to Mohammad's shape: introduction, background,
% protocol (argument only), experiments (everything measured), discussion,
% conclusion. This file is the opening and the setup; the result subsections
% follow in 04_below.tex until that is mined and replaced.
%
% Numbers here come from exp_AB/results.csv (the seven networks),
% exp_B/results.csv (the seventeen), exp_AA (LFR), exp_AC (SBM/dcSBM),
% exp_AD (Metis on real graphs). All verified against CSVs, not summaries.
% STYLE: no em-dashes, no AI-characteristic phrasing (see STYLE.md).
% ============================================================

Everything below runs under the protocol of Section~\ref{sec:protocol}. The
experiments answer three questions in order. What does the accounting cost, that
is, how much of the reported benefit in this literature is a property of how it
is measured. Where the pipeline is used in practice, does it help. And where it
does help, what are the conditions.

\subsection{Setup}\label{sec:experiments:setup}

\paragraph{Networks}
The quality measurements use the seven real networks of Table~\ref{tab:networks},
chosen to span an order of magnitude in average degree while remaining tractable
under a protocol that reruns every configuration with multiple seeds and computes
several controls per cell. Three carry reference communities: email-Eu-core has
department labels, and com-DBLP and com-Amazon carry the overlapping communities
distributed with them. The mechanism measurements, which require only one
partition per graph and no restarts, run on a larger suite of seventeen networks
extending to 25 million edges.
All graphs are reduced to their largest connected component and treated as
simple and undirected.

\begin{table}[t]
\centering
\caption{The seven networks carrying the quality measurements, ordered by average
degree. Sizes are of the largest connected component. GT marks the three with
reference communities.}
\label{tab:networks}
\small
\setlength{\tabcolsep}{5pt}
\begin{tabular}{@{}lrrrc@{}}
\toprule
Network & $n$ & $m$ & $d_{\mathrm{avg}}$ & GT\\
\midrule
com-Amazon    & 334{,}863 &   925{,}872 &  5.53 & $\bullet$\\
ca-HepTh      &   8{,}638 &    24{,}806 &  5.74 & \\
com-DBLP      & 317{,}080 & 1{,}049{,}866 &  6.62 & $\bullet$\\
ca-CondMat    &  21{,}363 &    91{,}286 &  8.55 & \\
email-Enron   &  33{,}696 &   180{,}811 & 10.73 & \\
wiki-Vote     &   7{,}066 &   100{,}736 & 28.51 & \\
email-Eu-core &      986 &    16{,}064 & 32.58 & $\bullet$\\
\bottomrule
\end{tabular}
\end{table}

Real networks cannot answer questions about where a transition lies, because
their density cannot be varied while their community structure is held fixed. Two
families of synthetic benchmark do that. LFR
graphs~\cite{lancichinetti2008benchmark} with $10^4$ vertices and degree exponent
$2.1$ are generated at nominal average degrees of $10$, $25$, $50$, $100$ and
$200$, three graphs per configuration, matching the generator settings of the
original synthetic sweep. Stochastic block models and degree-corrected block
models are generated at the same degrees, with degree and mixing calibrated per
cell to the realized values of the LFR graphs, in two variants each so that
degree heterogeneity and community-size heterogeneity form a $2 \times 2$
factorial. Realized degree, realized mixing and degree coefficient of variation
are recorded per graph and all results are reported against realized values.

\paragraph{Sparsifiers}
Seven methods are tested, chosen so that a negative result cannot be attributed
to the selection. Two are the methods whose quality claims motivate the pipeline:
local sparsification, L-Spar~\cite{satuluri2011local}, which lets each vertex keep
the $\lceil d_v^{\,e} \rceil$ incident edges of highest neighbourhood Jaccard
similarity, and the degree-based sampler DSpar~\cite{liu2023dspar}. Three more
are the methods that the largest independent benchmark of sparsification ranks
highest at preserving clustering structure, namely K-Neighbor, Local Degree and
Local Similarity~\cite{chen2024demystifying, hamann2016structure}, so that the
arms include the strongest candidates chosen by someone other than us. L-Spar and
Local Similarity are distinct methods despite the similar names, and L-Spar
always refers to the method of Satuluri et al. The remaining two are a
connectivity-preserving backbone, a spanning structure topped up to the retention
target either at random or by Jaccard similarity, which cannot disconnect a graph
by construction and is included as an instrument rather than as a competitor.
Uniform random deletion supplies the no-information baseline wherever a selection
rule has to be shown to be doing work.

Two classes are excluded and the exclusions bound what follows. Node-removing
methods such as $k$-core peeling are not tested, because they score their
partition on a smaller and denser vertex set, which is a different comparison
requiring the coverage control of Section~\ref{sec:protocol:retention}. Learned
sparsifiers~\cite{wu2020gsgan} are not tested, because matching them on realized
retention requires retraining at each operating point.

\paragraph{Detection algorithms}
Four algorithms choose their own number of communities:
Leiden~\cite{traag2019louvain} and Louvain~\cite{blondel2008fast}, which maximize
\eqref{eq:modularity}, Infomap~\cite{rosvall2008maps}, which minimizes a
description length, and label propagation~\cite{raghavan2007near}. One takes the
number as an input: Metis~\cite{karypis1998metis}, through the \texttt{pymetis}
bindings, which also imposes a balance constraint on part sizes. Where Metis is
used, $k$ is fixed to the same value in the sparsified and unsparsified arms, so
the two partitions have identical community counts by construction and no
granularity effect is possible.

On real networks $k$ has to be chosen, and we report both defensible choices: the
community count a modularity optimizer returns on the original graph, and, on the
three networks with reference communities, the number of those communities. The
two differ by as much as a factor of sixteen and are reported separately
throughout.

Of the algorithms behind the founding accuracy claim, Metis is tested here on
every network. For Graclus and MLR-MCL we could find no maintained implementation
that we were able to run, so MLR-MCL, the algorithm that claim is largest about,
remains untested by us. This is the most consequential gap in the paper and
Section~\ref{sec:discussion} returns to it.

\paragraph{Operating points}
Sparsifiers are compared at matched realized retention, with targets of $0.5$ and
$0.2$. Four of the seven methods cannot reach the lower target on a sparse graph,
because each vertex retains at least one incident edge, so a third operating
point is added per network at the largest floor among the arms, ensuring one
aggressive point at which all seven are comparable. Table~\ref{tab:floors} gives
the floors, and they are a result in their own right rather than a nuisance
parameter, since they bound how far the pipeline can be pushed on the graphs
where it would need to pay off most.

\begin{table}[t]
\centering
\caption{Hard retention floors, as a fraction of edges. Below these values the
method cannot run at all, whatever target it is given. Networks are ordered by
average degree, which governs the floor. DSpar and uniform deletion have no
floor.}
\label{tab:floors}
\small
\setlength{\tabcolsep}{4pt}
\begin{tabular}{@{}lrrrrr@{}}
\toprule
Network & $d_{\mathrm{avg}}$ & Backbone & L-Spar & Loc.\ Deg. & K-Neigh.\\
\midrule
com-Amazon    &  5.53 & 0.362 & 0.301 & 0.352 & 0.330\\
ca-HepTh      &  5.74 & 0.348 & 0.276 & 0.347 & 0.317\\
com-DBLP      &  6.62 & 0.302 & 0.244 & 0.301 & 0.277\\
ca-CondMat    &  8.55 & 0.234 & 0.186 & 0.234 & 0.218\\
email-Enron   & 10.73 & 0.186 & 0.159 & 0.186 & 0.179\\
wiki-Vote     & 28.51 & 0.070 & 0.067 & 0.070 & 0.069\\
email-Eu-core & 32.58 & 0.061 & 0.053 & 0.061 & 0.060\\
\bottomrule
\end{tabular}
\end{table}

\paragraph{Replication}
Stochastic sparsifiers are run with independent sparsification seeds and
stochastic detectors with independent detection seeds; deterministic methods are
run once and reused. Baselines receive the two comparisons of
Section~\ref{sec:protocol:compute}, the compute-matched restarts and the best of
a fixed number. Where Metis is used the comparison is the worst sparsified run
against the best baseline run over ten partitioner seeds, which is the strictest
form available and the one used throughout for that detector.

\paragraph{Pre-registration}
Each experiment has a design document written before its code, fixing what is
measured, what is predicted, and what outcome would end that line of work.
Several predictions were ours and failed, and where a result below contradicts a
prediction we made, it is reported as such.
